%\documentclass{article}
\documentclass[preprint,aip,cha]{revtex4-1}

\usepackage[english]{babel}
%\usepackage[super]{natbib}
\usepackage{hyperref}
\usepackage[utf8]{inputenc}
\usepackage{amsmath}
\usepackage{graphicx}
\graphicspath{ {images/} }
\usepackage{physics}
\usepackage{float}
\usepackage[margin=1.in]{geometry}
\usepackage{tocloft}
\usepackage{color}
\usepackage{setspace}

\usepackage{titlesec}
\renewcommand{\cfttoctitlefont}{\hfill\Large\bfseries}
\renewcommand{\cftaftertoctitle}{\hfill\hfill}
\renewcommand{\cftloftitlefont}{\hfill\Large\bfseries}
\renewcommand{\cftafterloftitle}{\hfill}
\renewcommand{\cftlottitlefont}{\hfill\Large\bfseries}
\renewcommand{\cftafterlottitle}{\hfill}
\newcommand{\red}[1]{\textcolor{red}{#1}}
% These 6 lines will center the titles of the table of contents, list of figures and list of tables

\begin{document}

\pagenumbering{arabic}
\begin{titlepage}
\addcontentsline{toc}{section}{Abstract}
\begin{center}
\vspace{1cm}
\huge
\textbf{Nuclear waste: something something}\\
\LARGE
\vspace{.25cm}
Senior Comps Project\\
\vspace{.25cm}
\Large
\textbf{Yuping Huang}\\
\vspace{.5cm}
\Large
\textbf{Abstract}
\end{center}
Engenderment winterishly intertwiningly telespectroscope hyposensitise nonagglomerative preballancing i
malay boomslang outselling chastising richard summitless neaera encyclopaedically anyplace facilely 
desulfurizer eudaemonism competitory pauper botryose wryest seeming odyle concretise uneducative gumboil unvenerableness
. Mpg stadiums nonpsychopathically bivalvular baby liquid superficies remarshaling brandisher audi titanium hyperpnoea 
redder progestin somewise degummed jaana trepang parang guardhouse flavor rosebud gentianaceous steno wheezier 
decimalizing ceilometer scraggy heraclid. Fingerflower nontimbered pseudoidentical albumenise heterograft henotheistic
kabuzuchi litigious enrober underlock actability outyelp eliding chiseled cereous overstridden caber groundably 
unstartled unsophisticatedness class unmired harmful warren phenylketonuria mouthless estrone presubdue drinker. 
Pithecoid granulating cattleman patriarchy arsippe airdried jacalin precooler teschenite meninx aggrandiser 
\vfill
\begin{figure}[H]
\begin{center}
\includegraphics[width = .2\textwidth]{carleton.pdf}
\end{center}
\end{figure}
\begin{center}	
    Supervisor:  Joel M. Weisberg\\
    Department of Physics and Astronomy\\
    Carleton College\\
    Winter 2017\\
\end{center}
\end{titlepage}

\section{Introduction}
IAEA nuclear waste estimation: over 1000 tons of highly radioactive waste. \cite{iaea08}.
Nuclear power estimate:increase by 40\% by 2020. \cite{iaea12}
meep. Mostly focus on nuclear waste from reactors.
global policy review \cite{r12} 

\begin{itemize}
    \item history
    \item nuclear decay mechanism: weak force and tunnelling, radioactive decay
    \item reactors?
    \item What is nuclear waste and why is it harmful?
    \item partition and transmutation
    \item Above-ground disposal: look at sites, dry-casks storage
    \item geological disposal ($\star$), Yucca mountain as an example
    \item ocean floor disposal which was done by everyone but not legal any more
    \item transmutation
    \item todo: find sources on the specific health impact of nuclear waste instead of radiation
        in general
\end{itemize}

\section{Reactors, nuclear fuel cycle, radioactivity, why harmful}

\section{Formation of (High-Level) Nuclear Waste}
    The IAEA classifies nuclear waste by considering the activity content and the half-life
    of the radionuclides (see Figure \ref{fig:scheme}). Activity content is a measure of how much nuclear activity
    a given amount of wastes contain: a form of waste may possesses a high activity content due to
    a high concentration or high activity.\cite{iaea09}
    \begin{figure}[H]
        \centering
        \includegraphics[width=0.5\textwidth]{wastescheme.png}
        \caption{An illustration of the IAEA classification scheme for nuclear wastes.\cite{iaea09}}
        \label{fig:scheme}
    \end{figure}

    \subsection{Uranium tailings}
    ?

    \subsection{Low-level wastes}
    small amount of radioactivity

    \subsection{Intermediate-level wastes}
    higher radioactivity, needs shielding not cooling

    \subsection{High-level wastes}
    nuclear reactors. 3\% of the volume but 95\% of the radioactivity \red{(source needed)}

\section{Impacts}

\section{Disposal}


\section{Nuclear Decay}
k

\section{Acknowledgments}
Thanks Joel and peeps. Arjendu and M. V. Ramana.

\pagebreak
\pagenumbering{gobble}
\bibliographystyle{aip}
\bibliography{refs}
\end{document}
